% answers_short.tex —— 2016 年·全国Ⅲ卷·文科 · 简短答案速查
% 由 tools/make_short_answers.py 自 answers.tex 的 \gksolhead/\gkanswer 提取，
% 并与 2016/2016_数学真题_全国3文_答案要点.md 逐题比对后生成，勿手改（改源 MD 后重跑）。

\begin{gknotice}
\textbf{使用说明}：本节只给最终答案便于快速核对，完整推导见解析版。答案为编写者独立推导，\textbf{非官方标准答案}。
\end{gknotice}

\begin{gknotice}
选择题答案中的 U 表示待核，未给出猜测答案，详情见解析版。
\end{gknotice}

\gkpart{选择题}
\begin{center}\begin{tabular}{ccccccccccccc}
\toprule
题号 & 1 & 2 & 3 & 4 & 5 & 6 & 7 & 8 & 9 & 10 & 11 & 12 \\
\midrule
答案 & C & D & A & D & C & D & A & B & D & B & B & U \\
\bottomrule
\end{tabular}\end{center}

\gkpart{填空题}
\begin{itemize}[leftmargin=2.6em]
  \item[\textbf{13.}] \(-10\)
  \item[\textbf{14.}] \(\dfrac\pi3\)
  \item[\textbf{15.}] \(4\)
  \item[\textbf{16.}] \(y=2x\)
\end{itemize}

\gkpart{解答题}
\begin{itemize}[leftmargin=2.6em]
  \item[\textbf{17.}] (1) \(a_2=\dfrac12,\ a_3=\dfrac14\)；(2) \(a_n=2^{1-n}\)
  \item[\textbf{18.}] (1) \(r\approx0.993\)，可用线性回归模型拟合；(2) \(\hat y=0.92+0.10t\)，预测约 \(1.82\) 亿吨
  \item[\textbf{19.}] (1) 见解析版；(2) \(\dfrac{4\sqrt5}{3}\)
  \item[\textbf{20.}] (1) 见解析版；(2) \(x=y^2\ (y\ne0)\) 或 \(x=y^2+1\)
  \item[\textbf{21.}] (1) 在 \((0,1)\) 上递增，在 \((1,+\infty)\) 上递减；(2) 见解析版；(3) 见解析版
\end{itemize}

\gkpart{选做题}
\begin{itemize}[leftmargin=2.6em]
  \item[\textbf{22.}] (1) \(60^\circ\)；(2) 见解析版（垂直平分线按通常定义存在的非退化构型）
  \item[\textbf{23.}] (1) \(C_1:\dfrac{x^2}{3}+y^2=1\)，\(C_2:x+y=4\)；(2) 最小值 \(\sqrt2\)，此时 \(P=(\dfrac32,\dfrac12)\)
  \item[\textbf{24.}] (1) \([-1,3]\)；(2) \(a\ge2\)
\end{itemize}
